\documentclass[12pt,a4paper]{article}
\usepackage[utf8]{inputenc}
\usepackage[T1]{fontenc}
\usepackage[portuguese]{babel}
\usepackage{geometry}
\geometry{a4paper, margin=2cm}
\usepackage{tikz}
\usetikzlibrary{shapes.geometric, shapes.multipart, arrows, positioning, fit, calc}
\usepackage{tabularx}
\usepackage{booktabs}
\usepackage{hyperref}
\usepackage{enumitem}

\title{\textbf{Definição do Processo de Desenvolvimento}}
\author{A1 de MC656}
\date{}

\begin{document}

\maketitle

\section{Integrantes do grupo}
\begin{itemize}
    \item \textbf{Giovanni Santos Scalabrin} - RA: 281210
    \item \textbf{Rodrigo Banin} - RA: 238257
    \item \textbf{Leandro Marcos Checchio de Souza} - RA: 244784
    \item \textbf{Lucca Maso Miranda} - RA: 281827
\end{itemize}

\section{Produto e contexto}
\subsection{Descrição do produto}
O sabrielGoares é um sistema de gerenciamento e auditoria de votações projetado para atender a diversos contextos institucionais, em alinhamento com o ODS 16 da ONU. Como cobre 3 cenários, abstraímos as lógicas de elegibilidade, peso dos votos, quóruns, sigilo, etc, o que facilita o uso para diferentes modelos eleitorais.

\subsection{Contexto de uso}
O sistema atende a três cenários com regras distintas:
\begin{enumerate}
    \item Eleição de Centro Acadêmico: Uma eleição com lista fechada, em que somente alunos podem votar. A votação ocorre por chapas, tem sigilo absoluto e a apuração usa maioria simples.
    \item Assembleia de Condomínio: Uma votação onde moradores e proprietários possuem votos com pesos proporcionais ao tamanho do imóvel. O processo inclui exige quóruns específicos para cada pauta, como a aprovação de obras. Além disso também possui votos por procuração.
    \item Orçamento Participativo Municipal: Um processo aberto a cidadãos locais, dividido em envio de propostas, análise de viabilidade técnica e votação de múltipla escolha. Os projetos escolhidos ficam sujeitos a um teto de orçamento. O voto é auditável, mas sem sigilo obrigatório.
\end{enumerate}

\subsection{Pressupostos}
\begin{itemize}
    \item Os eleitores e administradores possuem acesso à internet.
    \item A auditoria não quebra o sigilo dos votos onde a lei exige anonimato.
    \item O programa estrutura o processo eleitoral de forma que permita acomodar as diferenças específicas de cada eleição.
\end{itemize}

\subsection{Representações Visuais Iniciais}

\subsubsection*{Representação 1: Diagrama de Domínio}
O diagrama demonstra como o sistema isola a entidade de eleição das regras específicas de cada cenário.

\vspace{0.5cm}
\begin{center}
\begin{tikzpicture}[
    node distance=2.5cm and 1cm,
    class/.style={
        rectangle split, rectangle split parts=2,
        draw, align=left, top color=white, bottom color=blue!5, thick,
        text width=4cm, rounded corners=2pt
    }
]
    \node[class] (election) {election \nodepart{two} + title: string \\ + start: date \\ + end: date \\ + type: string \\ \vspace{1mm} + calculateresults()};
    \node[class, right=1.5cm of election] (vote) {vote \nodepart{two} + payload: string \\ + weight: float \\ + timestamp: date};
    \node[class, right=1.5cm of vote] (voter) {voter \nodepart{two} + id: string \\ + credentials: string \\ \vspace{1mm} + authenticate()};
    
    \node[class, below=2cm of election] (context) {$\langle\langle$interface$\rangle\rangle$ \\ contextrule \nodepart{two} + validateeligibility() \\ + calculateweight() \\ + checkquorum()};
    \node[class, right=1.5cm of context] (condo) {condorule \nodepart{two} + applyfractionalweight() \\ + allowproxy()};

    \draw[diamond-, thick] (election) -- node[above] {1..*} (vote);
    \draw[-, thick] (vote) -- node[above] {n..1} (voter);
    \draw[diamond-, thick] (election) -- node[left] {1..1} (context);
    \draw[-latex, dashed, thick] (condo) -- (context);
\end{tikzpicture}
\end{center}
\vspace{0.5cm}

\subsubsection*{Representação 2: Jornada do Usuário no Orçamento Participativo}
O diagrama ilustra o fluxo em que o cidadão envia uma proposta e, após a validação do administrador, retorna para votar.

\vspace{0.5cm}
\begin{center}
\begin{tikzpicture}[
    node distance=1.5cm and 1.5cm,
    stepnode/.style={rectangle, draw, rounded corners=4pt, align=center, fill=green!5, thick, text width=2.8cm, minimum height=1.5cm},
    arrow/.style={-latex, thick, auto}
]
    \node[stepnode] (s1) {cidadão envia proposta};
    \node[stepnode, right=of s1] (s2) {administrador faz análise técnica};
    \node[stepnode, right=of s2] (s3) {cidadão se autentica};
    \node[stepnode, below=of s3] (s4) {cidadão vota no teto};
    \node[stepnode, left=of s4] (s5) {sistema emite confirmação e recibo};

    \draw[arrow] (s1) -- (s2);
    \draw[arrow] (s2) -- (s3);
    \draw[arrow] (s3) -- (s4);
    \draw[arrow] (s4) -- (s5);
\end{tikzpicture}
\end{center}
\vspace{0.5cm}

\subsubsection*{Representação 3: Diagrama de Casos de Uso}
O diagrama mapeia as permissões do eleitor e do administrador no sistema.

\vspace{0.5cm}
\begin{center}
\begin{tikzpicture}[
    actor/.style={circle, draw, align=center, fill=gray!10, thick, text width=1.8cm, inner sep=2pt},
    usecase/.style={ellipse, draw, align=center, fill=yellow!10, thick, text width=3.5cm, inner sep=4pt},
    node distance=1.5cm and 3cm
]
    \node[actor] (eleitor) {eleitor};
    
    \node[usecase, right=3cm of eleitor] (votar) {registrar voto};
    \node[usecase, above=of votar] (submeter) {enviar proposta};
    \node[usecase, above=of submeter] (auth) {autenticar};
    
    \node[actor, right=3cm of votar] (admin) {admin};
    \node[usecase, below=of votar] (config) {configurar regras};
    \node[usecase, below=of config] (apurar) {apurar resultados};

    \draw[-, thick] (eleitor) -- (votar);
    \draw[-, thick] (eleitor) -- (submeter);
    \draw[-, thick] (eleitor) -- (auth);
    

    \draw[-, thick] (admin) -- (config);
    \draw[-, thick] (admin) -- (apurar);
    \draw[-, thick] (admin) -- (auth);
\end{tikzpicture}
\end{center}
\vspace{0.5cm}

\section{Processo de desenvolvimento proposto}

\subsection{Abordagem adotada}
A equipe usará o Kanban integrado ao Behavior-Driven Development (BDD). Como o contexto eleitoral envolve regras que variam entre os cenários, a abordagem precisa ser flexível, nesse sentido, o BDD resolve o problema da validação ao testar as regras de negócio diretamente no código, e o Kanban permite que os desenvolvedores assumam as tarefas no ritmo necessário.

\subsection{Relação entre atividades e fluxo}
O fluxo de trabalho organiza-se nas seguintes etapas:
\begin{enumerate}
    \item O Product Owner (PO) insere as demandas no backlog e semanalmente avalia se é preciso repriorizar, com o feedback do ambiente de Staging.
    \item A área de negócios atua na tradução dos requisitos, escrevendo os cenários de teste em Gherkin.
    \item Os desenvolvedores puxam essas funcionalidades do backlog e escrevem o código.
    \item O servidor de integração contínua executa os testes automaticamente. Se falhar, o PR fica bloqueado e a tarefa deve ser corrigida.
    \item Após a revisão, o CD leva o código ao ambiente de Staging.
    \item O PO analisa a funcionalidade em Staging e, se tudo estiver correto, autoriza a liberação para Produção.
    \item A equipe organiza uma retrospectiva a cada 15 dias para avaliar as métricas e ajustar o fluxo de trabalho se for preciso.
\end{enumerate}

\subsection{Justificativa}
Escrever os cenários em BDD antes de implementar evita erros nas regras de votação. O Kanban se alinha a isso pois foca em uma demanda por vez e reduz o peso do planejamento inicial. A automação de testes e deploy permite incluir novos cenários eleitorais sem afetar as regras existentes.

\section{Aplicação dos conceitos de Processo de Software}

\textbf{Tabela 1: Caracterização das atividades} \\
\begin{tabularx}{\textwidth}{|l|l|X|X|}
\hline
Atividade & Classificação & Objetivo & Quando/condição \\
\hline
especificar bdd & Fundamental & Traduzir as regras em testes em Gherkin. & Ao refinar o Backlog. \\
\hline
Implementar Código & Fundamental & Implementar para cumprir o BDD. & BDD aprovado no To-Do. \\
\hline
Revisar Código & Guarda-chuva & Garantir qualidade e segurança. & Ao abrir PR. \\
\hline
Validar Integração & Guarda-chuva & Executar testes no código. & Na pipeline de CI. \\
\hline
Aprovar Staging & Fundamental & Verificar se a interface atende às necessidades do usuário. & Disponível em Staging. \\
\hline
retrospectiva & Guarda-chuva & Melhorar o fluxo e remover gargalos. & A cada 15 dias. \\
\hline
\end{tabularx}

\vspace{0.5cm}

\textbf{Tabela 2: Execução das atividades} \\
\begin{tabularx}{\textwidth}{|X|X|X|X|X|}
\hline
Atividade & Participantes & Recursos & Entradas & Resultados \\
\hline
especificar bdd & PO & Google, documentos de regras & Regras do Domínio & Cenário BDD \\
\hline
Implementar Código & Desenvolvedores & IDE, Git & Cenário BDD & Código, Commits \\
\hline
Revisar Código & Dev Revisor & GitHub/GitLab & PR & Aprovação \\
\hline
Aprovar Staging & Product Owner & Amb. Staging & Pacote Staging & Release \\
\hline
retrospectiva & Toda a equipe & Quadro Kanban e Métricas & Gargalos observados & Ações de melhoria no Backlog \\
\hline
\end{tabularx}

\vspace{0.5cm}

\textbf{Tabela 3: Artefatos do processo} \\
\begin{tabularx}{\textwidth}{|l|X|X|}
\hline
Nome do artefato & Finalidade & Atividades em que é criado, atualizado ou utilizado \\
\hline
Especificação BDD & Definir as regras de votação em formato testável. & Criado e atualizado na especificação do BDD. Utilizado na implementação do código. \\
\hline
PR & Enviar alterações de código para revisão e integração. & Criado na implementação do código. Utilizado na revisão de código e na validação da integração. \\
\hline
Pacote de Release & Preparar a versão testada do sistema para deploy em produção. & Criado após a validação da integração. Utilizado na aprovação em Staging e deploy para prod. \\
\hline
\end{tabularx}

\section{Representação visual do processo}

\vspace{0.5cm}
\begin{center}
\begin{tikzpicture}[
    node distance=1.5cm and 1cm,
    process/.style={rectangle, draw, rounded corners, align=center, fill=blue!5, thick, text width=3cm, minimum height=1cm},
    decision/.style={diamond, draw, align=center, fill=yellow!5, thick, aspect=1.5, text width=2cm},
    arrow/.style={-latex, thick}
]
    \node[process] (bdd) {especificar bdd};
    \node[process, below=of bdd] (impl) {implementar\\funcionalidade};
    \node[process, below=of impl] (pr) {abrir pr};
    \node[process, below left=1cm and -0.5cm of pr] (ci) {validar\\integração};
    \node[process, below right=1cm and -0.5cm of pr] (rev) {revisar\\código};
    \node[decision, below=3cm of pr] (check) {aprovado?};
    \node[process, below=1.5cm of check] (cd) {deploy staging};
    \node[process, right=1.5cm of cd] (prod) {deploy produção};
    \node[process, right=1cm of bdd, fill=green!10] (retro) {retrospectiva};

    \draw[arrow] (bdd) -- (impl);
    \draw[arrow] (impl) -- (pr);
    \draw[arrow] (pr) -| (ci);
    \draw[arrow] (pr) -| (rev);
    \draw[arrow] (ci) |- (check);
    \draw[arrow] (rev) |- (check);
    \draw[arrow] (check) -- node[right] {sim} (cd);
    \draw[arrow] (check) -- node[above right] {não} ++(4,0) |- (impl);
    \draw[arrow] (cd) -- node[above, font=\scriptsize] {po valida} (prod);
    \draw[arrow, dashed] (prod) -| (retro);
    \draw[arrow, dashed] (retro) -- node[above, font=\scriptsize] {ajustes} (bdd);
\end{tikzpicture}
\end{center}
\vspace{0.5cm}

\section{Declaração de uso de IA Genereativa no processo de elaboração}
Utilizamos IA generativa para ajudar a comparar os métodos/abordagem de desenvolvimento, além disso usamos para ajuda na sintaxe do LaTeX para gerar os (lindos!) diagramas acima. De qualquer forma, os autores assumem \underline{total responsabilidade} pelo conteúdo deste trabalho.

\end{document}
